% !Mode:: "TeX:UTF-8"
% !TEX program  = xelatex
\begin{中文摘要}{LoRA 推理，适配器调度，预测驱动缓存，指数滑动平均，资源管理}
随着大语言模型的广泛应用，基于 PEFT（Parameter-Efficient Fine-Tuning）的技术，如 LoRA（Low-Rank Adaptation），被广泛用于支持多任务与个性化推理。在实际部署中，一个推理系统往往需要同时服务大量 LoRA 适配器，而 GPU 显存限制使得系统无法同时加载全部适配器，从而需要在运行过程中频繁进行适配器的加载与替换。该过程会引入显著的性能开销，成为多 LoRA 推理系统的关键瓶颈。

现有方法通常采用被动的适配器管理策略，例如基于 LRU 或 LFU 的缓存机制。这类方法仅依赖历史访问信息进行决策，缺乏对未来请求的建模能力，在动态或突发负载场景下容易导致频繁的缓存失效与适配器切换。

为了解决上述问题，本文提出一种基于预测的适配器调度方法。该方法利用指数滑动平均（EMA）对适配器访问趋势进行建模，并结合请求队列信息与适配器活跃状态构建统一的效用函数，同时引入显存感知机制对不同适配器进行归一化处理。在显存约束下，系统通过求解近似的效用最大化问题动态维护适配器驻留集合，从而优先保留高价值适配器。
\end{中文摘要}

\begin{英文摘要}{LoRA inference, adapter scheduling, prediction-driven caching, exponential moving average, resource management}
With the rapid adoption of large language models (LLMs), parameter-efficient fine-tuning (PEFT) methods such as LoRA (Low-Rank Adaptation) have been widely used to support multi-task and personalized inference. In practical deployment scenarios, a serving system often needs to manage a large number of LoRA adapters simultaneously. However, due to limited GPU memory, it is impossible to keep all adapters resident on GPU at the same time, resulting in frequent adapter loading and replacement during runtime. This process introduces significant overhead and becomes a major performance bottleneck in multi-LoRA serving systems.

Existing approaches typically adopt passive adapter management strategies, such as LRU- or LFU-based caching mechanisms. These methods rely only on historical access information and lack the ability to model future requests, making them prone to frequent cache misses and adapter switching under dynamic or bursty workloads.

To address these challenges, this paper proposes a prediction-based adapter scheduling method. The proposed approach utilizes Exponential Moving Average (EMA) to model adapter access trends, and further incorporates request queue information and adapter activity states into a unified utility function. In addition, a memory-aware normalization mechanism is introduced to handle heterogeneous adapter sizes. Under GPU memory constraints, the system dynamically maintains the resident adapter set by approximately solving a utility maximization problem, thereby prioritizing high-value adapters and reducing unnecessary adapter reloads.
\end{英文摘要}
\cleardoublepage
